\documentclass{article}
\usepackage{amssymb,amsmath}
\usepackage{graphicx}
\usepackage{enumerate}
\usepackage{amssymb,amsmath,amsthm}
\usepackage{graphicx}
\usepackage{tikz}
\usepackage{pgfplots}
\usepackage[utf8]{inputenc}
\usepackage[brazil]{babel}
\usepackage{enumerate}
\usepackage{color}
\usepackage{mathabx}
\usepackage{calc}
\usepackage{fullpage}

\usepackage{dsfont} % includes bb 1
\newcommand*\1{\mathds{1}}
\usepackage[brazil]{babel}

\newtheorem{theorem}{Teorema}[section]
\newtheorem{corollary}[theorem]{Corolário}
\newtheorem{lemma}[theorem]{Lema}
\newtheorem{proposition}[theorem]{Proposição}
\newtheorem{definition}[theorem]{Definição}
\newtheorem{notation}[theorem]{Notação}

\DeclareMathOperator{\Var}{Var}
\def\d{\mathrm{d}}

\begin{document}

\title{Probabilidade I}
\author{Segunda Prova}

\maketitle

\noindent Observa\c{c}\~oes:
\begin{itemize}
\item A prova terá duração de $3$ horas e meia.
\item Questões deverão ser respondidas de maneira matematicamente rigorosa e clara.
\item Todos os resultados provados em aula (exceto exercícios) poderão ser utilizados, salvo quando a questão pedir que algum destes seja provado. Nesse caso, todos os resultados anteriores poderão ser utilizados.
\item Caso uma questão seja composta de vários itens, cada um poderá ser resolvido independentemente, supondo a validade dos anteriores.
\end{itemize}

\vspace{4mm}
\noindent {\bf 1)} Sejam $X_1, X_2, \dots$ variáveis aleatórias independentes com distribuição Poisson de parâmetros $\lambda_1, \lambda_2, \dots$.
Definindo $S_n = X_1 + \dots + X_n$, mostre que
\begin{equation}
  \frac{S_n - E(S_n)}{\sqrt{\Var(S_n)}} \Rightarrow \mathcal{N}(0, 1)
\end{equation}
se e somente se $\sum_n \lambda_n = \infty$.

\vspace{2mm}
\noindent
Dica: Lembre que $P[\text{Poisson}(\lambda) = k] = \lambda^k e^{-\lambda}/k!$ e que se $X, Y$ são independentes e com distribuições Poisson($\lambda_1$) e Poisson($\lambda_2$) respectivamente, então $X + Y$ tem distribuição Poisson($\lambda_1 + \lambda_2$).

\vspace{4mm}
\noindent {\bf 2)} Sejam $X_1, X_2, \dots$ variáveis aleatórias i.i.d. e defina
\begin{equation}
  S_n = \sum_{i = 1}^n X_i, \text{ para $n \geq 1$}.
\end{equation}
\begin{enumerate}[\quad a)]
\item (Propriedade de Markov) Mostre que para quaisquer $k, n \geq 1$ e qualquer função $f: \mathbb{R}^k \to \mathbb{R}_+$ limitada,
  \begin{equation}
    E \big[ f(S_{n + 1}, \dots, S_{n + k}) | S_1, \dots, S_n \big]
    = E \big[ f(S_{n + 1}, \dots, S_{n + k}) | S_n \big].
  \end{equation}
\item (Martingal com respeito à filtração canônica) Mostre que se $X_i \in \mathcal{L}^1$ então, para qualquer $n \geq 1$,
  \begin{equation}
    E \big[ S_{n + 1} | S_n \big] = S_n.
  \end{equation}
\end{enumerate}

\pagebreak

\vspace{4mm}
\noindent {\bf 3)} Considere $f: \mathbb{R} \to \mathbb{R}_+$ contínua e de suporte compacto, que determine uma densidade na reta, ou seja, tal que $\int f(x) dx = 1$.
Gostaríamos de estimar o valor de $f$ em um ponto $x$ fixo, tendo acesso a apenas a amostras $X_1, \dots, X_N$ i.i.d. com distribuição $f(x) dx$.

Para isso, fixamos uma função $k: \mathbb{R} \to \mathbb{R}_+$ contínua e de suporte compacto, tal que $\int k(u) du = 1$ e, fixado $\theta > 0$, definimos o estimador
\begin{equation}
  F_\theta^N(x) := \frac{1}{N} \sum_{i = 1}^N \frac{1}{\theta}
  k \Big( \frac{x - X_i}{\theta} \Big),
\end{equation}
que é uma variável aleatória.
Nosso objetivo é estimar, para $x$ fixado,
\begin{equation}
  P \big[ \big| F_\theta^N(x) - f(x) \big| \geq t \big].
\end{equation}
\begin{enumerate}[\quad a)]
\item Mostre que
  \begin{equation}
    P \big[ \big| F_\theta^N(x) - f(x) \big| \geq t \big]
    \leq \frac{E \big( ( F_\theta^N(x) - f(x) )^2 \big)}{t^2}.
  \end{equation}
\item Mostre que
  \begin{equation}
    E \big( ( F_\theta^N(x) - f(x) )^2 \big)
    = \big( E (F_\theta^N(x) - f(x)) \big)^2
    + \Var(F_\theta^N).
  \end{equation}
\item Lembrando que $k$ é contínua e de suporte compacto e supondo ainda que $\int u k(u) du = 0$ e que $f \in C^1$, mostre que
  \begin{equation}
    \big| E (F_\theta^N(x) - f(x)) \big|
    \leq \frac{\theta^2}{2} \lVert f'' \rVert_\infty \int u^2 k(u) du.
  \end{equation}
  Dica: desenvolva a série de Taylor para $f$ em torno de $x$.
\item Ainda assumindo que $k$ é contínua e de suporte compacto, mostre que
  \begin{equation}
    \Var(F_\theta^N(x)) \leq \frac{1}{N \theta^2}
    \Var \Big( k \Big( \frac{x - X_1}{\theta} \Big) \Big)
    \leq \frac{1}{N \theta^2}
    E \Big( k \Big( \frac{x - X_1}{\theta} \Big)^2 \Big)
    \leq \frac{\lVert f \rVert_\infty}{N \theta} \int k^2(u) du.
  \end{equation}
\end{enumerate}


\end{document}
